\subsection{Calibration}

Calibration is a cornerstone of reliable measurements. Instruments used for data collection should always be either **calibrated** or **verified against a trusted, calibrated reference**.

In practice, not all laboratory instruments are externally calibrated, as accredited traceable calibration can be expensive. For example, calibrating a common handheld multimeter can cost more than the instrument itself, while calibration of a vector network analyzer (VNA) can easily range from thousands to tens of thousands of euros, often performed annually.

This underscores the principle that expensive instruments must be handled with care: even a minor accident, such as dropping a VNA calibration kit, can necessitate recalibration, resulting in significant cost.

It is therefore the experimenter’s responsibility to ensure that instruments are either properly calibrated or verified against a known reference. This ensures that the measured data reflect the physical quantity of interest, rather than artifacts arising from inaccurate or malfunctioning instruments.
